\section{The constants and the order of their choice}\label{sec:constants}

Proposition~\ref{prop:U1} asserts the existence of one choice of constants and one construction with all the stated properties. Table~\ref{tab:constant-choices} describes that choice; its numerical restrictions are properties of this construction. Each row may use only preceding rows, and stage-indexed choices are made for every fixed~$j\ge4$ before the entry levels are chosen.

\refstepcounter{table}\label{tab:constant-choices}
\begingroup
\begin{tabbing}
\hspace{2.4em}\=\hspace{0.27\linewidth}\=\kill
\textbf{No.}\>\textbf{Choice}\>\textbf{Dependence, restrictions and use}\\[0.4em]
1\>\parbox[t]{0.26\linewidth}{\raggedright $G\subset U(n_G)$, $dh$, metric, $\dimG,\casimir,b_0,\slope$; $\mathsf b>0$}\>\parbox[t]{0.67\linewidth}{\raggedright The prescribed realization and penalty; $\slope>0$. These are fixed throughout.}\\[0.45em]
2\>\parbox[t]{0.26\linewidth}{\raggedright Kernel cutoffs $\delta_{\rm fr},\varepsilon_{\rm pred}$ and axial factor $a_{\rm ax}$}\>\parbox[t]{0.67\linewidth}{\raggedright $G$; positive sufficiently small cutoffs and $a_{\rm ax}\ge16$, with the branch conditions at every patch. The penalty remains fixed. Proposition~\ref{prop:U1}, B1 of Theorem~\ref{thm:B}.}\\[0.45em]
3\>\parbox[t]{0.26\linewidth}{\raggedright Homogeneous label bases, norms, local sections, stencils}\>\parbox[t]{0.67\linewidth}{\raggedright $G$ and degree; Taylor matching, reality and $H_4$ covariance. These precede the ratios needed for their bounds. Propositions~\ref{prop:U3}--\ref{prop:U4} and B1.}\\[0.45em]
4\>\parbox[t]{0.26\linewidth}{\raggedright $\mathfrak q_j,\varpi_j,\sigma_j$, clipping constants, $R_j^{\rm rel},w_j^{\rm path}$, $A_j^{\rm path},C_j,\omega_j^{\rm RG}$}\>\parbox[t]{0.67\linewidth}{\raggedright Rows 1--3 and $j$; $\mathfrak q_j\ge6$, $\varpi_j\ge2$, both tending to infinity, $\sigma_j=100$, $C_j\ge1$, $R_j^{\rm rel}w_j^{\rm path}\ge12$; positive clipping constants, relaxation lengths, widths and path bounds. Choose $\omega_j^{\rm RG}\ge0$ above the one-loop residual. All constants entering ratio bounds are fixed here, including those for the lattice estimates of Propositions~\ref{prop:P6}, \ref{prop:P10}, \ref{prop:P12} and~\ref{prop:P13}.}\\[0.45em]
5\>\parbox[t]{0.26\linewidth}{\raggedright Odd nondecreasing stage ratios $\ell_j^{\rm st}$}\>\parbox[t]{0.67\linewidth}{\raggedright Rows 1--4; exceed every ratio threshold and satisfy $\omega_j^{\rm RG}/\log\ell_j^{\rm st}\le\min(\slope/4,j^{-2})$ and the logarithmic bound on ratio constants in the level-data definition. Propositions~\ref{prop:U1}--\ref{prop:U2}.}\\[0.45em]
6\>\parbox[t]{0.26\linewidth}{\raggedright Antitone exponents $\theta_j$; entry bounds}\>\parbox[t]{0.67\linewidth}{\raggedright Rows 1--5; take the finite maximum of the remaining coupling and smallness thresholds at each stage. Further entry bounds may be raised. Propositions~\ref{prop:U1}--\ref{prop:U4}, \ref{prop:P6}, \ref{prop:P10}, \ref{prop:P12}, \ref{prop:P13}.}\\[0.45em]
7\>\parbox[t]{0.26\linewidth}{\raggedright $\bar\beta$, entry levels, minimum periods}\>\parbox[t]{0.67\linewidth}{\raggedright Rows 1--6; ensure $\refbeta_m>1$, $t_m\le1$ and every entry condition. Entry levels can be increased after earlier constants are fixed. Proposition~\ref{prop:U2}.}\\[0.45em]
8\>\parbox[t]{0.26\linewidth}{\raggedright $J_m,\ell_m,D_m,a_m,L_p$, $\reflevel{d}$}\>\parbox[t]{0.67\linewidth}{\raggedright Rows 1--7; gentle growth, patch and cover fits, and volume divisibility. Require $0<\theta_{J_m}\le J_m^{-1}$, $J_{\reflevel{d}}\ge d+3$ and sufficient derivative orders thereafter. Propositions~\ref{prop:U2}--\ref{prop:U4}.}\\[0.45em]
9\>\parbox[t]{0.26\linewidth}{\raggedright $\refbeta_m,\cutoffsize_m,t_m,B_m$, $\epsilon_m^{\rm reg},\epsilon_m^{\rm for}$, $\Lambda_m^{\rm loop},\mathfrak B_m$}\>\parbox[t]{0.67\linewidth}{\raggedright The exact level-data formulas determine these quantities from rows 1--8. There is no new choice. Propositions~\ref{prop:U1}--\ref{prop:U3}.}\\[0.45em]
10\>\parbox[t]{0.26\linewidth}{\raggedright Path actions, classes, recursion, trajectories, $\beta_{N,p}$, $w_{\rm tun}$}\>\parbox[t]{0.67\linewidth}{\raggedright Rows 1--9 and Propositions~\ref{prop:U1}--\ref{prop:U2}; the common level-zero endpoint and the specified recursion apply for every finite $N,p$. Used in Propositions~\ref{prop:U3}--\ref{prop:U4} and B1.}\\[0.45em]
11\>\parbox[t]{0.26\linewidth}{\raggedright Marked transfers, terminal maps, renormalized insertions}\>\parbox[t]{0.67\linewidth}{\raggedright Rows 1--10; normalize at the fixed $\reflevel{d}$. The entire Wilson family is fixed before a volume sequence is prescribed. Used in every continuum conclusion.}\\[0.45em]
12\>\parbox[t]{0.26\linewidth}{\raggedright $\mathsf p$, then $\mathsf q,S,F_d,\partial_a,\mathcal L$}\>\parbox[t]{0.67\linewidth}{\raggedright Prescribe arbitrary strictly increasing $\mathsf p$ after row 11. Proposition~\ref{prop:P5} chooses one refinement and table, using the full inner limit of Proposition~\ref{prop:U4}.}\\[0.45em]
13\>\parbox[t]{0.26\linewidth}{\raggedright Comparison radii, frames, $J$, charges, $c_\star,R_d,T$}\>\parbox[t]{0.67\linewidth}{\raggedright Row 12 and the fixed family; Proposition~\ref{prop:P6}, then Proposition~\ref{prop:P7}. Comparison radii may be refined while the volume refinement remains fixed.}\\[0.45em]
14\>\parbox[t]{0.26\linewidth}{\raggedright Slab operators, $\phases,\nu,S^\xi$, then $\xi_*$}\>\parbox[t]{0.67\linewidth}{\raggedright Row 13 and Proposition~\ref{prop:P8}; select $\xi_*$ only for a one-dimensional OS ground space.}\\[0.45em]
15\>\parbox[t]{0.26\linewidth}{\raggedright All-phase theories and continuations; $\Theta,\tracespectrum$}\>\parbox[t]{0.67\linewidth}{\raggedright Rows 12--14; Proposition~\ref{prop:P9}, then Proposition~\ref{prop:P10}. Centering uses the same selected phase. Also used in Proposition~\ref{prop:P11}.}\\[0.45em]
16\>\parbox[t]{0.26\linewidth}{\raggedright One OPE prescription and coefficients; contacts, real bases, positive scales $\delta_q^{\rm ct}$}\>\parbox[t]{0.67\linewidth}{\raggedright Rows 12--14 and the fixed family; Proposition~\ref{prop:P12}, then Proposition~\ref{prop:P13} with unchanged projections. Require $4\delta_{q+1}^{\rm ct}<\delta_q^{\rm ct}$ and $2(q+1)\delta_q^{\rm ct}\le\delta_q^{\rm pres}$.}\\[0.45em]
17\>\parbox[t]{0.26\linewidth}{\raggedright Thresholds $N_0,k_0,r_0$; estimate constants; local distribution seminorms}\>\parbox[t]{0.67\linewidth}{\raggedright Exactly the fixed data preceding each threshold in its statement. There is no dependence on a later position, volume or state whenever uniformity is asserted. These are chosen for individual estimates.}\\
\end{tabbing}
\endgroup

The order is acyclic because the constants entering a ratio bound are fixed before that ratio, whereas the remaining smallness requirements are met by increasing entry levels. At each fixed stage only finitely many additional bounds are combined, and increasing entry levels preserves earlier inequalities. Proposition~\ref{prop:U1} makes these choices jointly with the represented densities, marked transfers and additional lattice estimates; the displayed numerical restrictions alone are not a sufficiency criterion for other constructions. Rows 1--11 precede the prescribed volume sequence in row 12; the subsequent choices in rows 12--16 are existential. Row 9 is determined by formulas, and row 17 consists of thresholds for fixed requests, rather than construction constants.
